\scanpage{343}
\mspage{82}
% archive sheet number: 342
vrai, mais fastidieux à expliciter.

Les sueurs froides sont terminées~!

En résumé, le formalisme de \inserted{$\underline{\mathrm{Ch}}_\Sigma(X)$ et de} $\underline{\mathrm{Ch}}_\Sigma(X;a,b)$ développé jusqu'à présent donne satisfaction complète sur les points suivants~:

a) Bon formalisme de composition des chemins (qui était dégueulasse avant)

b) Traitement élégant de la construction des catégories de fractions $\deleted{\unsure{X}}{\unsure{M|X}}\Sigma^{-1}$. Mais je m'attends \unsure{que} ce sera moins joli qu'avant, quand il s'agira de remplacer un $(\deleted{X}M,\Sigma)$ par des
\[
\underset{\substack{\shortparallel\\\underline{\mathrm{Hom}}(X,M)}}{M(X)},\Sigma_X
\]
($X$ une catégorie d'indices). C'est pour y pallier qu'il \inserted{me} faudra développer des énoncés de comparaison, entre la catégorie $\underline{\mathrm{Ch}}_\Sigma$ \unsure{construite|construit} ici, et d'autres \unsure{\textsuperscript{$q$}} dont le type dans $\mathrm{Hot}$ se prête aux calculs de façon plus directe, comme dans l'état premier de la question.

De plus, on sera sûrement obligé d'en venir à des Ind-catégories, pour pouvoir \unsure{sauver|suivre} la \unsure{variance} de $M(X)\Sigma_X^{-1}$ pour $X$ variable.

c) Bon formalisme des ``homotopies'' et des compositions d'homotopies (ce qui est une autre façon de parler de chemins et de leur composition).
